\documentclass[11pt,a4paper]{article}
\input{../../sablony/pripravy/priprava.sty}
\usepackage{tikz,pgfplots}
\pgfplotsset{compat=1.18}
\nastavPripravu{1. ročník}{Kvadratické rovnice a nerovnice}{107}{Viètovy vztahy}
\begin{document}
\blok{Viètovy vztahy}
Součet a součin kořenů; rychlá kontrola a jednoduché určování kořenů bez vzorce.
\begin{definiceBox}
\textbf{Viètovy vztahy:} pro kořeny kvadratické rovnice $ax^2+bx+c=0$ ($a\ne0$) platí $x_1+x_2=-\frac ba$ a $x_1x_2=\frac ca$.
\end{definiceBox}
\blok{Odvození z již známého vzorce pro kořeny}
Předpokládejme $D\ge0$. Ze vzorce pro kořeny máme
\[x_1=\frac{-b+\sqrt D}{2a},\qquad x_2=\frac{-b-\sqrt D}{2a}.\]
Potom součet a součin získáme přímo výpočtem:
\[
\begin{aligned}
x_1+x_2&=\frac{-b+\sqrt D-b-\sqrt D}{2a}=-\frac ba,\\
x_1x_2&=\frac{b^2-D}{4a^2}
=\frac{b^2-(b^2-4ac)}{4a^2}=\frac ca.
\end{aligned}
\]
Zatím nepotřebujeme rozklad trojčlenu -- ten si odvodíme až v další hodině.
\blok{Metoda a její omezení}
Viètovy vztahy jsou výborné pro rychlou zkoušku a hledání jednoduchých kořenů.
Nejsou však náhradou za diskriminant, pokud se kořeny nedají odhadnout.
\begin{prikladReseniBox}[title={Ukázka -- zkontroluj oba kořeny bez dosazování}]
Pro rovnici $x^2-5x+6=0$ ověřte navržené kořeny $2$ a $3$ pomocí Viètových vztahů.
\tcblower
\[\begin{aligned}
x_1+x_2&=2+3=5=-\frac{-5}{1},\\
x_1x_2&=2\cdot3=6=\frac61,\\
K&=\{2;3\}.
\end{aligned}\]
\end{prikladReseniBox}
\blok{Společně řešené příklady}
\begin{prikladReseniBox}[title={Příklad 1}]
$x^{2} - 5x + 6=0$
\tcblower
\[
\begin{aligned}
x_1+x_2&=5,\quad x_1x_2=6\\[2pt]
x&=2\ \vee\ x=3\\[2pt]
K&=\{2;3\}
\end{aligned}
\]
\end{prikladReseniBox}
\begin{prikladReseniBox}[title={Příklad 2}]
$2x^{2} - 6x - 20=0$
\tcblower
\[
\begin{aligned}
x_1+x_2&=3,\quad x_1x_2=-10\\[2pt]
x&=-2\ \vee\ x=5\\[2pt]
K&=\{-2;5\}
\end{aligned}
\]
\end{prikladReseniBox}
\begin{prikladReseniBox}[title={Příklad 3}]
$x^{2} + 5x + 4=0$
\tcblower
\[
\begin{aligned}
x_1+x_2&=-5,\quad x_1x_2=4\\[2pt]
x&=-4\ \vee\ x=-1\\[2pt]
K&=\{-4;-1\}
\end{aligned}
\]
\end{prikladReseniBox}
\blok{Samostatné procvičování}
\begin{zadaniReseniBox}[title={Úloha 1}]
$x^{2} - 10x + 24=0$
\tcblower
\[
\begin{aligned}
x_1+x_2&=10,\quad x_1x_2=24\\[2pt]
x&=4\ \vee\ x=6\\[2pt]
K&=\{4;6\}
\end{aligned}
\]
\end{zadaniReseniBox}
\begin{zadaniReseniBox}[title={Úloha 2}]
$3x^{2} - 15x + 12=0$
\tcblower
\[
\begin{aligned}
x_1+x_2&=5,\quad x_1x_2=4\\[2pt]
x&=1\ \vee\ x=4\\[2pt]
K&=\{1;4\}
\end{aligned}
\]
\end{zadaniReseniBox}
\blok{Volitelná rezerva}
\poznamka{Studenti si vymyslí obdobný příklad se změnou jednoho koeficientu a vysvětlí, jak se změní počet kořenů či řešení.}
\blok{Domácí úkol}
\begin{zadaniReseniBox}[title={Domácí příprava}]
\begin{enumerate}[label=\alph*),itemsep=2pt]
\item $x^{2} - 11x + 24=0$
\item $2x^{2} - 14x=0$
\item $x^{2} + 7x + 10=0$
\item $3x^{2} - 3x - 6=0$
\item $x^{2} - 2x + 1=0$
\item $2x^{2} + 14x + 24=0$
\end{enumerate}
\tcblower
\begin{enumerate}[label=\alph*),itemsep=2pt]
\item
\[\begin{aligned}
&\left(x - 8\right) \cdot \left(x - 3\right)=0\\[2pt]
x&=3\ \vee\ x=8\\[2pt]
K&=\{3;8\}
\end{aligned}\]
\item
\[\begin{aligned}
&2x \cdot \left(x - 7\right)=0\\[2pt]
x&=0\ \vee\ x=7\\[2pt]
K&=\{0;7\}
\end{aligned}\]
\item
\[\begin{aligned}
&\left(x + 2\right) \cdot \left(x + 5\right)=0\\[2pt]
x&=-5\ \vee\ x=-2\\[2pt]
K&=\{-5;-2\}
\end{aligned}\]
\item
\[\begin{aligned}
&\left(x - 2\right) \cdot \left(3x + 3\right)=0\\[2pt]
x&=-1\ \vee\ x=2\\[2pt]
K&=\{-1;2\}
\end{aligned}\]
\item
\[\begin{aligned}
&\left(x - 1\right)^{2}=0\\[2pt]
x&=1\ \vee\ x=1\\[2pt]
K&=\{1\}
\end{aligned}\]
\item
\[\begin{aligned}
&\left(x + 4\right) \cdot \left(2x + 6\right)=0\\[2pt]
x&=-3\ \vee\ x=-4\\[2pt]
K&=\{-4;-3\}
\end{aligned}\]
\end{enumerate}
\end{zadaniReseniBox}
\poznamka{Orientační čas: úvod 5 min, společné příklady 18 min, samostatně 15 min, kontrola a shrnutí 7 min. Rezerva je dobrovolná. Číselné parametry modelových aplikací nejsou skutečně naměřená data.}
\end{document}
